\chapter{Electronic Effects and Reactive Intermediates}\label{ch:b1:electronic-effects}

Trifluoroethanoic acid, \ce{CF3COOH}, is about ten thousand times
stronger an acid than ethanoic acid, \ce{CH3COOH}, although the acidic
O--H bond is three bonds away from the fluorine atoms. Phenol is a
million times more acidic than ethanol, and an ammonia molecule bound to
a benzene ring --- aniline --- is a million times less basic than one
bound to a methyl group. Organic chemistry is full of such contrasts, and
two effects explain most of them: the pull of electronegative atoms along
the $\sigma$ bonds, and the spreading of electron pairs over $\pi$
systems. The same two effects decide which bonds break, which
intermediates form and where reagents attack: they are the grammar of the
\omterm{def:b1:elementary-steps:elementary-step}{reaction mechanisms} of the next chapters.

\begin{recall}
\omterm{def:b1:periodicity:electronegativity}{Electronegativity} and its trends (\cref{ch:b1:periodicity}); \omterm{def:b1:lewis-resonance-vsepr:lewis}{Lewis
structures}, \omterm{def:b1:lewis-resonance-vsepr:formal-charge}{formal charges} and resonance (\cref{ch:b1:lewis-resonance-vsepr});
$\mathrm{p}K_a$ and the strength of acids and bases
(\cref{ch:b1:predominant-reaction}). The school volume (grade 12) drew
\omterm{def:b1:electronic-effects:curly-arrow}{curly arrows} from a \omterm{def:b1:electronic-effects:nucleophile}{nucleophile} to an \omterm{def:b1:electronic-effects:nucleophile}{electrophile}.
\end{recall}

\section{Inductive effects}

\begin{definition}[Inductive effect]\label{def:b1:electronic-effects:inductive}
The \emph{inductive effect}\index{inductive effect} is the shift of
electron density along $\sigma$ bonds towards an electronegative atom or
group. An attracting group (F, Cl, O, \ce{NO2}, \ce{CF3}) has a $-I$
effect; alkyl groups and negatively charged atoms push electrons, a $+I$
effect. The effect weakens quickly with the number of bonds in between.
\end{definition}

\begin{proposition}[Inductive effect and acidity]\label{prop:b1:electronic-effects:inductive-pka}
An attracting group near the acidic site of a carboxylic acid
stabilises the carboxylate ion by spreading its negative charge, and
lowers the $\mathrm{p}K_a$: the more attracting groups, the closer, the
stronger the acid.
\end{proposition}

\begin{proof}
$\mathrm{p}K_a$ measures the position of $\mathrm{RCOOH} \rightleftharpoons
\mathrm{RCOO^-} + \ce{H+}$. A group that draws electron density away from
the carboxylate spreads its charge over more atoms, which lowers its
energy more than that of the neutral acid; the equilibrium moves to the
right and $K_a$ increases. The measured series confirms it.
\end{proof}

\begin{omfigure}
\begin{tikzpicture}
\begin{axis}[
  xbar, width=0.78\linewidth, height=5.2cm, xmin=0, xmax=5.4,
  symbolic y coords={TFA,TCA,DCA,MCA,AA}, ytick=data,
  yticklabels={\ce{CF3COOH}, \ce{CCl3COOH}, \ce{CHCl2COOH}, \ce{CH2ClCOOH}, \ce{CH3COOH}},
  xlabel={$\mathrm{p}K_a$ at \qty{25}{\celsius}}, bar width=9pt,
  nodes near coords, nodes near coords align={horizontal},
  every node near coord/.append style={font=\scriptsize},
  tick label style={font=\small}, label style={font=\small}, enlarge y limits=0.15]
\addplot[fill=omDef!45, draw=omDef] coordinates {(0.52,TFA) (0.51,TCA) (1.35,DCA) (2.87,MCA) (4.76,AA)};
\end{axis}
\end{tikzpicture}

{\small The halogenated ethanoic acids. Each chlorine lowers the
$\mathrm{p}K_a$; with three, the acid is about ten thousand times
stronger than ethanoic acid. Trifluoro- and trichloroethanoic acids
are equally strong within the uncertainty of measurements on such \omterm{def:b1:predominant-reaction:strong-acid}{strong
acids}.}
% ledger: pka:ethanoic, pka:chloroethanoic, pka:dichloroethanoic, pka:trichloroethanoic, pka:trifluoroethanoic
\end{omfigure}

Alkyl groups work the other way, a little: propanoic acid ($\mathrm{p}K_a
= 4.89$) is slightly weaker than ethanoic acid (4.76), and methanoic acid
(3.74), with no alkyl group, is stronger than both.
% ledger: pka:propanoic, pka:methanoic

\section{Mesomeric effects}

When an atom bearing a \omterm{def:b1:lewis-resonance-vsepr:lewis}{lone pair}, or a charge, is bound to a $\pi$
system, its electrons can be shared with that system; \omterm{def:b1:lewis-resonance-vsepr:resonance}{resonance
structures} (\cref{ch:b1:lewis-resonance-vsepr}) describe the sharing.

\begin{definition}[Mesomeric effect, donor and acceptor groups]\label{def:b1:electronic-effects:mesomeric}
The \emph{mesomeric effect}\index{mesomeric effect} is the shift of
electron density through a conjugated $\pi$ system, shown by \omterm{def:b1:lewis-resonance-vsepr:resonance}{resonance
structures}. A \emph{donor group}\index{donor group} gives a \omterm{def:b1:lewis-resonance-vsepr:lewis}{lone pair} to
the system, a $+M$ effect (\ce{-OH}, \ce{-O^-}, \ce{-NH2}, \ce{-OR}); an
\emph{acceptor group}\index{acceptor group} draws electrons from it through
a $\pi$ bond to an electronegative atom, a $-M$ effect (\ce{-NO2},
\ce{-C=O}, \ce{-C#N}).
\end{definition}

\begin{omfigure}
\setchemfig{atom sep=2em}
\schemestart
\chemfig{CH_3-C(=[1]O)-[7]O^{-}}
\arrow{<->}
\chemfig{CH_3-C(-[1]O^{-})=[7]O}
\schemestop
\hspace{2.5em}
\schemestart
\chemfig{*6(=-=(-O^{-})-=-)}
\arrow{<->}
\chemfig{*6((=O)-=-\chemabove{}{\scriptstyle\ominus}-=-)}
\schemestop

\medskip
\schemestart
\chemfig{H_3C-C(=[1]O)-[7]NH_2}
\arrow{<->}
\chemfig{H_3C-C(-[1]O^{-})=[7]NH_2^{+}}
\schemestop

{\small \omterm{def:b1:lewis-resonance-vsepr:resonance}{Resonance structures} of the ethanoate ion (the charge shared by
two equivalent oxygens), of the phenoxide ion (the charge spread into the
ring; one of the three ring structures shown), and of an amide (the
nitrogen \omterm{def:b1:lewis-resonance-vsepr:lewis}{lone pair} shared with the C=O, which makes amides poor bases).}
\end{omfigure}

\begin{proposition}[Mesomeric effect and acidity]\label{prop:b1:electronic-effects:mesomeric-pka}
An acid whose conjugate base spreads its charge by resonance is stronger
than one whose base cannot: carboxylic acids ($\mathrm{p}K_a$ about 4--5)
are far stronger than alcohols (about 16), and phenol (10.0) lies between.
An \omterm{def:b1:electronic-effects:mesomeric}{acceptor group} in the position \textit{ortho} or \textit{para} to the OH
of a phenol strengthens it more than in the position \textit{meta}.
\end{proposition}

\begin{proof}
In an alkoxide the charge sits on one oxygen; in a carboxylate it is
shared by two oxygens; in a phenoxide it is shared by the oxygen and three
ring carbons, less electronegative than oxygen, a smaller gain. A
\ce{NO2} group at the \textit{ortho} or \textit{para} carbon can take the
charge itself through a \omterm{def:b1:lewis-resonance-vsepr:resonance}{resonance structure} with \ce{C=N+(-O^-)-O^-}; at
\textit{meta}, no \omterm{def:b1:lewis-resonance-vsepr:resonance}{resonance structure} brings the charge onto the carbon
that carries \ce{NO2}, and only the \omterm{def:b1:electronic-effects:inductive}{inductive effect} remains. The measured
values: 2-nitrophenol 7.23, 4-nitrophenol 7.16, 3-nitrophenol 8.36,
against 9.99 for phenol and 15.9 for ethanol.
\end{proof}
% ledger: pka:ethanol, pka:phenol, pka:2-nitrophenol, pka:3-nitrophenol, pka:4-nitrophenol

The same reasoning applies to bases. Methylamine ($\mathrm{p}K_a$ of its
ammonium 10.66) is a stronger base than ammonia (9.25, \cref{ch:b1:predominant-reaction}),
the methyl group pushing electrons onto nitrogen; in aniline the \omterm{def:b1:lewis-resonance-vsepr:lewis}{lone pair}
is shared with the ring and the ammonium is far more acidic (4.60).
Pyridine (5.23) keeps its \omterm{def:b1:lewis-resonance-vsepr:lewis}{lone pair}, but on a nitrogen held in a planar
ring with more \textit{s} character, less available than in an amine.
% ledger: pka:methylammonium, pka:anilinium, pka:pyridinium, dfg:NH4+_ao

\section{Breaking a bond: reactive intermediates}

\begin{definition}[Homolysis and heterolysis]\label{def:b1:electronic-effects:cleavage}
A \omterm{def:b1:lewis-resonance-vsepr:lewis}{covalent bond} breaks by \emph{homolysis}\index{homolysis} when each atom
keeps one of the two bonding electrons, giving two \omterm{def:b1:electronic-effects:intermediates}{radicals}, and by
\emph{heterolysis}\index{heterolysis} when one atom keeps both, giving a
cation and an anion.
\end{definition}

\begin{definition}[Reactive intermediates]\label{def:b1:electronic-effects:intermediates}
A \emph{reactive intermediate}\index{reactive intermediate} is a species
formed in one step of a mechanism and consumed in a later one, too
reactive to accumulate (\cref{ch:b1:elementary-steps}). A
\emph{carbocation}\index{carbocation} has a carbon with three bonds and a
positive charge (six \omterm{def:b1:quantum-numbers:configuration}{valence electrons}); a
\emph{carbanion}\index{carbanion} a carbon with three bonds, a \omterm{def:b1:lewis-resonance-vsepr:lewis}{lone pair}
and a negative charge; a \emph{radical}\index{radical} an atom with an
\omterm{def:b1:quantum-numbers:unpaired}{unpaired electron}.
\end{definition}

\begin{definition}[Class of a carbon]\label{def:b1:electronic-effects:carbon-class}
A carbon bound to one other carbon is a \emph{primary
carbon}\index{primary carbon}, to two a \emph{secondary
carbon}\index{secondary carbon}, to three a \emph{tertiary
carbon}\index{tertiary carbon}. A \omterm{def:b1:electronic-effects:intermediates}{carbocation}, a \omterm{def:b1:electronic-effects:intermediates}{radical} or a
halogenoalkane takes the class of the carbon concerned.
\end{definition}

\begin{omfigure}
\begin{tikzpicture}[font=\small]
  % carbocation: planar, empty p orbital
  \begin{scope}
    \fill[omProp!15, draw=omProp] (0,0.15) ellipse (0.18 and 0.75);
    \fill[omProp!15, draw=omProp] (0,-0.15) ellipse (0.18 and 0.75);
    \draw[thick] (0,0) -- (-1.0,-0.25) node[left] {R};
    \draw[thick] (0,0) -- (0.95,-0.35) node[right] {R};
    \draw[thick] (0,0) -- (0.35,0.45) node[above right] {R};
    \node[fill=white, inner sep=1pt] at (0,0) {\ce{C+}};
    \node at (0,-1.75) {carbocation: planar,};
    \node at (0,-2.15) {empty \textit{p} orbital};
  \end{scope}
  % radical: nearly planar, one electron
  \begin{scope}[xshift=4.6cm]
    \draw[omMeth] (0,0.15) ellipse (0.18 and 0.75);
    \draw[omMeth] (0,-0.15) ellipse (0.18 and 0.75);
    \fill (0,0.55) circle (0.05);
    \draw[thick] (0,0) -- (-1.0,-0.25) node[left] {R};
    \draw[thick] (0,0) -- (0.95,-0.35) node[right] {R};
    \draw[thick] (0,0) -- (0.35,0.45) node[above right] {R};
    \node[fill=white, inner sep=1pt] at (0,0) {C};
    \node at (0,-1.75) {radical: almost planar,};
    \node at (0,-2.15) {one electron};
  \end{scope}
  % carbanion: pyramidal with a lone pair
  \begin{scope}[xshift=9.2cm]
    \draw[thick] (0,0) -- (-0.95,-0.55) node[left] {R};
    \draw[thick] (0,0) -- (0.95,-0.55) node[right] {R};
    \draw[thick] (0,0) -- (0.2,-0.9) node[below] {R};
    \fill[omDef!15, draw=omDef] (0,0.55) ellipse (0.22 and 0.45);
    \fill (-0.07,0.65) circle (0.045); \fill (0.07,0.65) circle (0.045);
    \node[fill=white, inner sep=1pt] at (0,0) {\ce{C-}};
    \node at (0,-1.75) {carbanion: pyramidal,};
    \node at (0,-2.15) {lone pair};
  \end{scope}
\end{tikzpicture}

{\small Geometries of the three \omterm{def:b1:electronic-effects:intermediates}{reactive intermediates} of carbon (R: H or
an alkyl group). The flat \omterm{def:b1:electronic-effects:intermediates}{carbocation} can be attacked on either face, a
point that matters for the stereochemistry of
\cref{ch:b1:nucleophilic-substitution}.}
\end{omfigure}

\begin{proposition}[Stability of carbocations]\label{prop:b1:electronic-effects:carbocation-order}
\omterm{def:b1:electronic-effects:intermediates}{Carbocations} are stabilised by groups that give them electrons: tertiary
$>$ secondary $>$ primary $>$ methyl; a \omterm{def:b1:electronic-effects:intermediates}{carbocation} next to a double bond
(allylic) or a benzene ring (benzylic) is further stabilised by
resonance, and one next to an atom with a \omterm{def:b1:lewis-resonance-vsepr:lewis}{lone pair} (\ce{R-O-CH2+}) even
more.
\end{proposition}

\begin{proof}
Each alkyl group pushes electron density towards the electron-poor
carbon ($+I$ effect, and a related interaction of its C--H bonds with the
empty orbital, described in the Year~2 volume). An allylic cation has
two \omterm{def:b1:lewis-resonance-vsepr:resonance}{resonance structures} putting the charge on either end carbon; a
benzylic cation has four, three in the ring; an oxygen \omterm{def:b1:lewis-resonance-vsepr:lewis}{lone pair} forms a
fourth bond to the cationic carbon, \ce{R-O+=CH2}, in which every atom has
an octet.
\end{proof}

\omterm{def:b1:electronic-effects:intermediates}{Carbanions} follow the opposite order (methyl $>$ primary $>$ secondary
$>$ tertiary) and are stabilised by \omterm{def:b1:electronic-effects:mesomeric}{acceptor groups}; \omterm{def:b1:electronic-effects:intermediates}{radicals} follow the
same order as \omterm{def:b1:electronic-effects:intermediates}{carbocations}, less markedly.

\section{Curly arrows, nucleophiles and electrophiles}

\begin{definition}[Curly arrows]\label{def:b1:electronic-effects:curly-arrow}
In a mechanism, a \emph{curly arrow}\index{curly arrow} shows the
movement of an electron pair: it starts on a \omterm{def:b1:lewis-resonance-vsepr:lewis}{lone pair} or a bond and ends
on an atom (where a \omterm{def:b1:lewis-resonance-vsepr:lewis}{lone pair} or a new bond forms) or between two atoms (a
new bond). A \emph{half-headed arrow}\index{half-headed arrow} shows the
movement of a single electron, in \omterm{def:b1:electronic-effects:intermediates}{radical} reactions.
\end{definition}

\begin{method}[Drawing curly arrows]\label{met:b1:electronic-effects:curly-arrows}
\begin{enumerate}
  \item Start from electrons: a \omterm{def:b1:lewis-resonance-vsepr:lewis}{lone pair}, a $\pi$ bond or a $\sigma$
        bond, never from a positive charge or an atom without electrons.
  \item End where the pair goes: a new bond between the two atoms, or a
        \omterm{def:b1:lewis-resonance-vsepr:lewis}{lone pair} on the more electronegative atom when a bond breaks.
  \item Count charges after each step: the atom giving a \omterm{def:b1:lewis-resonance-vsepr:lewis}{lone pair} to a
        new bond gains $+1$, the atom receiving a pair as a \omterm{def:b1:lewis-resonance-vsepr:lewis}{lone pair}
        gains $-1$; the total charge is conserved.
  \item Never exceed an octet on C, N, O: when a pair arrives, another
        must leave.
\end{enumerate}
\end{method}

\begin{omfigure}
\setchemfig{atom sep=2.2em}
\schemestart
\chemfig{H_3C-C(-[2]H)(-[6]H)-[@{b}]@{br}Br}
\arrow{->[heterolysis]}
\chemfig{H_3C-C(-[2]H)(-[6]H)^{+}}\+\chemfig{Br^{-}}
\schemestop
\chemmove{\draw[->, shorten <=2pt, shorten >=2pt](b) .. controls +(90:6mm) and +(90:6mm) .. (br);}

\medskip
\schemestart
\chemfig{H@{o}O^{-}}\+\chemfig{@{h}H-@{oc}O-H}
\arrow{<=>}
\chemfig{H_2O}\+\chemfig{HO^{-}}
\schemestop
\chemmove{\draw[->, thick, shorten <=2pt, shorten >=2pt](o) .. controls +(-80:8mm) and +(-100:8mm) .. (h);}
\par\vspace{5mm}

{\small \omterm{def:b1:electronic-effects:curly-arrow}{Curly arrows}. Top: \omterm{def:b1:electronic-effects:cleavage}{heterolysis} of the C--Br bond, the pair going
to bromine. Bottom: a proton transfer from water to hydroxide, written as
the attack of a \omterm{def:b1:lewis-resonance-vsepr:lewis}{lone pair} of the base on the hydrogen (the bond O--H of
water then breaks, its pair staying on oxygen; the second arrow is left
for \cref{exo:b1:electronic-effects:4}).}
\end{omfigure}

\begin{definition}[Nucleophile, electrophile, leaving group]\label{def:b1:electronic-effects:nucleophile}
A \emph{nucleophile}\index{nucleophile} is a species that gives an
electron pair to an atom other than hydrogen to form a bond (a \omterm{def:b1:lewis-resonance-vsepr:lewis-acid}{Lewis
base}); an \emph{electrophile}\index{electrophile} is a species that
accepts such a pair (a \omterm{def:b1:lewis-resonance-vsepr:lewis-acid}{Lewis acid}, often a carbon bearing a partial
positive charge). The \emph{nucleophilicity}\index{nucleophilicity} of a
species is its reactivity as a nucleophile, measured by \omterm{def:b1:rate-laws:order}{rate constants}.
A \emph{leaving group}\index{leaving group} is the group that departs
with the \omterm{def:b1:lewis-resonance-vsepr:lewis}{bonding pair} when a bond to carbon breaks heterolytically.
\end{definition}

\begin{proposition}[Nucleophilicity and basicity]\label{prop:b1:electronic-effects:nucleophilicity-basicity}
For \omterm{def:b1:electronic-effects:nucleophile}{nucleophiles} that attack through the same atom, \omterm{def:b1:electronic-effects:nucleophile}{nucleophilicity}
follows basicity: \ce{CH3O-} $>$ \ce{OH-} $>$ \ce{CH3COO-} $>$ \ce{H2O}. It
fails across a column of the periodic table in \omterm{def:b1:intermolecular-forces-solvents:solvent-classes}{protic solvents}, where
large, polarisable, weakly solvated \omterm{def:b1:electronic-effects:nucleophile}{nucleophiles} react faster although they
are weaker bases (\ce{I-} $>$ \ce{Br-} $>$ \ce{Cl-} $>$ \ce{F-}), and for
bulky bases, which are poor \omterm{def:b1:electronic-effects:nucleophile}{nucleophiles}.
\end{proposition}

\begin{proof}
Basicity measures the equilibrium of the attack on a proton,
\omterm{def:b1:electronic-effects:nucleophile}{nucleophilicity} the rate of attack on a carbon; both increase with the
availability of the electron pair, hence the parallel within one atom.
The rate also depends on the energy needed to strip the solvent from the
\omterm{def:b1:electronic-effects:nucleophile}{nucleophile} and on the distortion of its electron cloud towards a
crowded carbon: small anions such as \ce{F-} are strongly bound to \omterm{def:b1:intermolecular-forces-solvents:solvent-classes}{protic
solvent} molecules (\cref{ch:b1:intermolecular-forces-solvents}), and
bulky groups hinder the approach to carbon but not to a proton.
\end{proof}

A good \omterm{def:b1:electronic-effects:nucleophile}{leaving group} is a \omterm{def:b1:predominant-reaction:strong-acid}{weak base}, stable once it carries the pair:
\ce{I-}, \ce{Br-}, \ce{Cl-}, water, sulfonates; \ce{OH-} and \ce{RO-}, \omterm{def:b1:predominant-reaction:strong-acid}{strong
bases}, are poor \omterm{def:b1:electronic-effects:nucleophile}{leaving groups} --- the reason why alcohols must be
activated before substitution (\cref{ch:b1:alcohol-activation}).

\section{Exercises}

\begin{exercise}[$\star$]\label{exo:b1:electronic-effects:1}
Classify each carbon of 2-methylbutane and of 2-bromo-2-methylpropane as
primary, secondary, tertiary or quaternary (bound to four carbons). Which
\omterm{def:b1:electronic-effects:intermediates}{carbocation} forms most easily from each bromide of formula \ce{C4H9Br}?
\end{exercise}

\begin{exercise}[$\star$]\label{exo:b1:electronic-effects:2}
From the $\mathrm{p}K_a$ values, compute the ratio of the \omterm{def:b1:predominant-reaction:acidity-constant}{acidity
constants} of chloroethanoic and ethanoic acids, and of trifluoroethanoic
and ethanoic acids.
\end{exercise}

\begin{exercise}[$\star$]\label{exo:b1:electronic-effects:3}
Draw the \omterm{def:b1:lewis-resonance-vsepr:resonance}{resonance structures} of the ethanoate ion, of the 4-nitrophenoxide
ion (with the charge on an oxygen of the nitro group) and of the allyl
cation \ce{CH2=CH-CH2+}.
\end{exercise}

\begin{exercise}[$\star$]\label{exo:b1:electronic-effects:4}
Complete the \omterm{def:b1:electronic-effects:curly-arrow}{curly arrows} of the proton transfer of the figure and of
\ce{NH3 + H3O+ -> NH4+ + H2O}. Check the charges.
\end{exercise}

\begin{exercise}[$\star\star$]\label{exo:b1:electronic-effects:5}
Rank by increasing acidity: ethanol, phenol, ethanoic acid,
4-nitrophenol, trichloroethanoic acid, water ($\mathrm{p}K_a$ of the
couple \ce{H2O}/\ce{OH-}, 14.0). Justify each step with an inductive or
mesomeric argument.
\end{exercise}

\begin{exercise}[$\star\star$]\label{exo:b1:electronic-effects:6}
Rank by increasing basicity: aniline, methylamine, ammonia, pyridine.
Explain the position of aniline with a \omterm{def:b1:lewis-resonance-vsepr:resonance}{resonance structure}.
\end{exercise}

\begin{exercise}[$\star\star$]\label{exo:b1:electronic-effects:7}
Why is 3-nitrophenol weaker than 4-nitrophenol but still stronger than
phenol? Draw the \omterm{def:b1:lewis-resonance-vsepr:resonance}{resonance structures} that justify both statements.
\end{exercise}

\begin{exercise}[$\star\star$]\label{exo:b1:electronic-effects:8}
Rank the \omterm{def:b1:electronic-effects:intermediates}{carbocations} \ce{CH3+}, \ce{(CH3)3C+}, \ce{CH2=CH-CH2+},
\ce{CH3CH2+}, \ce{C6H5CH2+}, \ce{CH3OCH2+}, and justify.
\end{exercise}

\begin{exercise}[$\star\star$]\label{exo:b1:electronic-effects:9}
Identify the \omterm{def:b1:electronic-effects:nucleophile}{nucleophile} and the \omterm{def:b1:electronic-effects:nucleophile}{electrophile} in: \ce{CH3Br + OH-};
\ce{H2O + H+}; ethanal with a cyanide ion; \ce{BF3 + NH3}. Draw the
\omterm{def:b1:electronic-effects:curly-arrow}{curly arrows}.
\end{exercise}

\begin{exercise}[$\star\star\star$]\label{exo:b1:electronic-effects:10}
A solution contains \qty{0.010}{mol/L} of trichloroethanoic acid. Is the
weak-acid formula $\mathrm{pH} = \frac12(\mathrm{p}K_a - \log C)$ valid?
Compute the pH properly and the \omterm{def:b1:predominant-reaction:dissociation}{degree of dissociation}.
\end{exercise}

\begin{exercise}[$\star\star\star$]\label{exo:b1:electronic-effects:11}
The $\mathrm{p}K_a$ of alcohols measured in water increase from methanol
(15.5) to ethanol (15.9), propan-2-ol (17.1) and 2-methylpropan-2-ol
(19.2). Give two explanations, one through the \omterm{def:b1:electronic-effects:inductive}{inductive effect} of the
alkyl groups on the alkoxide, one through the \omterm{def:b1:intermolecular-forces-solvents:dissolution}{solvation} of the alkoxide.
% ledger: pka:methanol, pka:ethanol, pka:2-propanol, pka:tBuOH
\end{exercise}

\begin{exercise}[$\star\star\star$]\label{exo:b1:electronic-effects:12}
Explain why amides are neither basic on nitrogen nor nucleophilic,
whereas amines are both, using the \omterm{def:b1:lewis-resonance-vsepr:resonance}{resonance structure} of the amide.
Which atom of an amide is protonated in \omterm{def:b1:predominant-reaction:strong-acid}{strong acid}?
\end{exercise}

\section{Problem: Why Trifluoroethanoic Acid Is So Strong}

\begin{problem}[{Weekend problem --- the inductive series of the
halogenated acids, carboxylates against alkoxides, the three
nitrophenols, a predominant-reaction calculation, and the ratio of the
acidity constants of trifluoroethanoic and ethanoic acids}]
\label{pb:b1:electronic-effects:1}
Data ($\mathrm{p}K_a$ at \qty{25}{\celsius}): ethanoic acid 4.76,
chloroethanoic 2.87, dichloroethanoic 1.35, trichloroethanoic 0.51,
trifluoroethanoic 0.52, ethanol 15.9, phenol 9.99, 2-nitrophenol 7.23,
3-nitrophenol 8.36, 4-nitrophenol 7.16; $\mathrm{p}K_e = 14.00$.

\medskip
\noindent\textbf{Part I --- The inductive series.}
\begin{enumerate}
  \item Compute the ratio of $K_a$ for each step from ethanoic to
        trichloroethanoic acid.
  \item Is the effect of each chlorine the same? Comment.
  \item Explain the trend with the \omterm{def:b1:electronic-effects:inductive}{inductive effect}.
  \item Why would a chlorine on the carbon of a longer chain farthest
        from COOH have almost no effect?
  \item Fluorine is more electronegative than chlorine. What does the
        comparison of the trifluoro- and trichloro- acids show, and why
        is it hard to measure such $\mathrm{p}K_a$ precisely?
  \item In which form is each of these acids at pH 3?
\end{enumerate}

\medskip
\noindent\textbf{Part II --- Carboxylates and alkoxides.}
\begin{enumerate}[resume]
  \item Draw the two \omterm{def:b1:lewis-resonance-vsepr:resonance}{resonance structures} of the ethanoate ion.
  \item What do they predict for the two C--O bond lengths of the ion?
  \item Why has the ethoxide ion no equivalent structures?
  \item Compute the ratio of the $K_a$ of ethanoic acid and ethanol.
  \item Place phenol between the two and justify with the phenoxide
        \omterm{def:b1:lewis-resonance-vsepr:resonance}{resonance structures}.
  \item Which of ethanol, phenol and ethanoic acid react almost totally
        with a hydrogencarbonate solution ($\mathrm{p}K_a$ 6.37 for
        \ce{CO2,H2O}/\ce{HCO3-})?
\end{enumerate}

\medskip
\noindent\textbf{Part III --- The nitrophenols.}
\begin{enumerate}[resume]
  \item Rank the three nitrophenols and phenol by acidity.
  \item Draw a \omterm{def:b1:lewis-resonance-vsepr:resonance}{resonance structure} of 4-nitrophenoxide with the charge on
        the nitro group.
  \item Show that no such structure exists for 3-nitrophenoxide.
  \item Why is 3-nitrophenol still more acidic than phenol?
  \item 2-Nitrophenol is slightly weaker than 4-nitrophenol although both
        have the \omterm{def:b1:electronic-effects:mesomeric}{mesomeric effect}; propose an explanation involving its
        OH and the nearby nitro group.
  \item At pH 7.5, which nitrophenols are mostly in the anion form? Why
        does that make 4-nitrophenol a \omterm{def:b1:titration-methods:indicator}{colour indicator}?
\end{enumerate}

\medskip
\noindent\textbf{Part IV --- Trifluoroethanoic acid in water.}
\begin{enumerate}[resume]
  \item Compute the pH of trifluoroethanoic acid at \qty{0.010}{mol/L}
        (do not assume weak dissociation).
  \item Compute the pH of ethanoic acid at \qty{0.010}{mol/L}.
  \item Write the reaction of trifluoroethanoic acid with sodium ethanoate
        and compute its constant.
  \item State the ratio $K_a(\ce{CF3COOH})/K_a(\ce{CH3COOH})$, as a power
        of ten and as a number.
\end{enumerate}
\end{problem}
